% =============================================================================
% Note on this combined document, Summary, Preface, Affirmation
% =============================================================================

\mysection{Note on this combined document}
This document combines two things produced at different points in the same thesis project.
Part I (Chapters 1--6 and Appendices A--B) is the master thesis mid-term report as submitted on
April 26, 2026, covering method development and preliminary results on the NIST CoCr AM XCT
benchmark dataset (\texttt{set1sample2raw}). Its content is reproduced essentially verbatim from that
submission. Part II (Chapters 7--13) is the PODFAM validation phase chapter, completed after the
mid-term and reporting full-volume results on real PODFAM process-monitoring XCT data, a separate
piece of work using a different pseudo-labelling method (Bernsen instead of two-stage Otsu), added
here as a continuation rather than a rewrite of Part I.

Three specific values in Part I were found, while preparing this combined document, to no longer
match the working codebase's actual configuration and have been corrected here rather than left
silently wrong: the NLM denoising filter-strength factor (Section~3.3.4, Appendix~B), the
validation/test split percentages (Section~3.2.2, Appendix~B), and the checkpoint output directory
name (Appendix~B). Each correction is marked in place. Appendix A's repository structure has also
been updated to reflect the codebase's current layout, which has been substantially restructured
(a \url{src/} package and a \url{scripts/} entry-point folder were introduced) since the
mid-term submission; the original Appendix A tree described an earlier, now-superseded layout.
Chapter 4's figures (4.1--4.12) are retained as captioned placeholders --- the specific NIST-phase
source images referenced by those captions are not present in the current repository state and have
not been regenerated for this document, since doing so was outside the scope of what was requested;
the PODFAM phase's own figures (Part II) are all freshly generated from the current pipeline and are
embedded in full. All other content in Part I, including its reported quantitative results, is
reproduced as originally submitted and was not re-verified against a pipeline re-run.

The bibliography at the end of Part I has been extended into a single, deduplicated reference list
covering both parts; in-text citations within Part II use the same shared \texttt{references.bib}
file as Part I, so every citation in this document resolves through one numbered list.

\clearpage
\mysection{Summary}
This master's thesis develops and validates an automated pipeline for the detection of internal
defects in metal additive manufacturing (AM) components using X-ray Computed Tomography (XCT) data.
The work addresses a critical industrial challenge: the reliable and scalable identification of gas
porosity, lack-of-fusion voids, and micro-cracks in high-resolution XCT volumes acquired from
components fabricated by laser powder bed fusion (LPBF).

The primary dataset is the NIST CoCr AM XCT benchmark dataset (\texttt{set1sample2raw}), comprising
900 slices at $1010 \times 980$ pixel resolution with a raw intensity range of $[0, 59631]$. No
expert-annotated ground truth masks are available for this dataset; defect detection labels are
therefore generated automatically using a two-stage Otsu thresholding strategy applied within the
metal component region only, followed by morphological cleaning and connected component filtering.

A four-stage preprocessing pipeline has been implemented and fully validated on the complete
900-slice dataset in 578.5 seconds on CPU. The pipeline comprises percentile normalisation,
polynomial beam hardening correction (degree 3), polar-coordinate ring artefact suppression (radius
15 pixels), and non-local means denoising ($h = 0.6 \times \hat{\sigma}$, adaptive per slice ---
corrected from the mid-term submission's stated value; see the Note on this combined document). The
pipeline achieves a noise standard deviation reduction from 20952 to 0.385, a factor of
$\times 54{,}500$, and an SNR improvement from 1.756 to 1.773. A systematic parameter sensitivity
analysis across NLM filter strength, BHC polynomial degree, and ring filter radius confirms the
selected configuration as the optimal operating point for this dataset.

Patch extraction ($256 \times 256$ pixels, 50\% overlap, stratified 1:3 foreground-to-background
sampling) and a six-transform data augmentation pipeline have been implemented and validated on the
NIST CoCr data, producing 15,840 patches at $64 \times 64$ and 3,860 patches at $128 \times 128$ from
20 sampled slices.

A 2D U-Net detection model has been implemented in PyTorch and initial training on the automatically
generated labels is underway. After 20 training epochs, the combined Dice-Focal loss formulation
achieves a validation Dice score of 0.68, compared to 0.41 for the standard binary cross-entropy
baseline and 0.38 for the Otsu label generator itself. This confirms that the U-Net is already
generalising beyond the capability of its own training labels, the fundamental goal of the
pseudo-label approach.

The complete pipeline is version-controlled on GitHub, tracked with MLflow, and produces thesis-ready
quantitative outputs including intensity histograms, SNR plots, parameter tuning comparisons, patch
extraction grids, and augmentation validation figures.

This mid-term summary is extended, in Part II of this document, by a completed validation phase on
real PODFAM metal AM XCT data: full-volume classical-threshold comparison across up to seven samples,
a 2D U-Net trained on Bernsen pseudo-labels with defect-content-ranked sample selection (best
validation Dice 0.600), held-out test evaluation, and an additional out-of-distribution inference
check that surfaced and diagnosed a reproducible Bernsen calibration failure mode.

\clearpage
\mysection{Preface}
This thesis is carried out as a degree project for the Master of Engineering Science with a
specialization in Robotics and Automation at the Department of Engineering Science, University West,
Sweden. The work was conducted within the PODFAM research project during the spring semester of
2026, under the supervision of Amit Kumar Mishra and examination by Yongcui Mi.

The project was initiated with the goal of developing an automated pipeline for defect detection in
additively manufactured (AM) metal components using X-ray computed tomography (XCT) data. For model
development, I used benchmark datasets from the National Institute of Standards and Technology (NIST
AM Bench), which provided diverse and well-annotated XCT scans suitable for training deep learning
models. To ensure industrial relevance and robustness, the trained models were subsequently tested on
PODFAM project datasets, consisting of XCT volumetric scans of AM metal parts collected within
ongoing process monitoring research at University West.

I would like to express my sincere gratitude to my supervisor, Prof. Amit Kumar Mishra and Prof.
Karthikeyan Thalavai Pandian for his continuous guidance, constructive feedback, and invaluable
support throughout the project. His expertise in machine learning and industrial automation has been
instrumental in shaping the direction of this work. I also extend my thanks to my examiner, Yongcui
Mi, for her insightful comments and encouragement during the course of the thesis.

I am grateful to the PODFAM project team and collaborators for their engagement and for providing
access to XCT data and domain expertise, which greatly enhanced the practical impact of this
research. Special thanks also go to the Department of Engineering Science at University West for
providing the resources, facilities, and academic environment that made this work possible.

Finally, I would like to thank my family and friends for their patience, encouragement, and
unwavering support throughout my studies.

\clearpage
\mysection{Affirmation}
This master degree report, \emph{Automated XCT Data Analysis for Defect Detection in Additively
Manufactured Metal Parts}, was written as part of the master degree work needed to obtain a Master of
Science with specialization in Robotics degree at University West. All material in this report that
is not my own is clearly identified and used in an appropriate and correct way. The main part of the
work included in this degree project has not previously been published or used for obtaining another
degree.

\vspace{2cm}
\noindent Job George Konnoth Joseph \hfill Date: April 26, 2026

\clearpage
